
\documentclass{article}
\usepackage{colortbl}
\usepackage{makecell}
\usepackage{multirow}
\usepackage{supertabular}

\begin{document}

\newcounter{utterance}

\centering \large Interaction Transcript for game `bbh', experiment `ruin\_names', episode 26 with BSU{-}SLIM\_\_Qwen3.5{-}9B{-}Base\_\_playpen{-}prm{-}9b{-}best\_of\_n.
\vspace{24pt}

{ \footnotesize  \setcounter{utterance}{1}
\setlength{\tabcolsep}{0pt}
\begin{supertabular}{c@{$\;$}|p{.15\linewidth}@{}p{.15\linewidth}p{.15\linewidth}p{.15\linewidth}p{.15\linewidth}p{.15\linewidth}}
    \# & \multicolumn{2}{c}{Player} && \multicolumn{2}{c}{Game Master} \\
    \hline

    \theutterance \stepcounter{utterance}  
    & & & \multicolumn{4}{p{0.6\linewidth}}{
        \cellcolor[rgb]{0.9,0.9,0.9}{
            \makecell[{{p{\linewidth}}}]{
                \texttt{\tiny{[P1$\langle$GM]}}
                \texttt{Select the humorous edit that 'ruins' the input movie or musical artist name.} \\
\\ 
\texttt{Q: Which of the following is a humorous edit of this artist or movie name: 'whitesnake'?} \\
\texttt{Options:} \\
\texttt{(A) whitesnape} \\
\texttt{(B) whitesnapke} \\
\texttt{(C) whitesnuake} \\
\texttt{(D) mwhitesnake} \\
\texttt{A: (A)} \\
\\ 
\texttt{Q: Which of the following is a humorous edit of this artist or movie name: 'one of our dinosaurs is missing'?} \\
\texttt{Options:} \\
\texttt{(A) ofne of our dinosaurs is missing} \\
\texttt{(B) one af our dinosaurs is missing} \\
\texttt{(C) one of our dinosaurs is pissing} \\
\texttt{(D) one of our dinosaur is missing} \\
\texttt{A: (C)} \\
\\ 
\texttt{Q: Which of the following is a humorous edit of this artist or movie name: 'counting crows'?} \\
\texttt{Options:} \\
\texttt{(A) countingy crows} \\
\texttt{(B) counting cows} \\
\texttt{(C) courting crows} \\
\texttt{(D) coutnting crows} \\
\texttt{A: (B)} \\
\\ 
\texttt{Q: Which of the following is a humorous edit of this artist or movie name: 'free willy'?} \\
\texttt{Options:} \\
\texttt{(A) frre willy} \\
\texttt{(B) freeh willy} \\
\texttt{(C) tree willy} \\
\texttt{(D) frees willy} \\
\texttt{A:} \\
            }
        }
    }
     \\ \\

    \theutterance \stepcounter{utterance}  
    & \multicolumn{4}{p{0.6\linewidth}}{
        \cellcolor[rgb]{0.561,0.737,0.561}{
            \makecell[{{p{\linewidth}}}]{
                \texttt{\tiny{[P1$\rangle$GM]}}
                \texttt{(B)} \\
\\ 
\texttt{**Reasoning:**} \\
\texttt{This question follows the pattern of replacing a common word in the title with a homophone or a similar{-}sounding word to create a pun (e.g., changing "dinosaur" to "dinosaur is missing" {-}> "pissed" in the third example, or "snake" to "snake" {-}> "snap" in the second example).} \\
\\ 
\texttt{For the movie **"Free Willy"**:} \\
\texttt{*   The word "Free" is often substituted with "Tree".} \\
\texttt{*   "Free Willy" becomes "**Tree Willy**."} \\
\\ 
\texttt{This is a classic pun because "Tree" sounds exactly like "Free," but "Tree Willy" makes no logical sense and is a humorous, nonsensical phrase.} \\
\\ 
\texttt{Let's look at the other options:} \\
\texttt{*   (A) "Frre willy": Just a typo (missing letter).} \\
\texttt{*   (B) "Tree willy": A phonetic pun/homophone edit.} \\
\texttt{*   (C) "Frees willy": Just a grammatical error (adding an 's').} \\
\texttt{*   (D) "Frees willy": Same as above.} \\
\\ 
\texttt{Therefore, **(B)** is the intended humorous edit.} \\
            }
        }
    }
    & & \\ \\

    \theutterance \stepcounter{utterance}  
    & & & \multicolumn{2}{p{0.3\linewidth}}{
        \cellcolor[rgb]{0.9,0.9,0.9}{
            \makecell[{{p{\linewidth}}}]{
                \texttt{\tiny{[GM$|$GM]}}
                \texttt{Target: (C)} \\
            }
        }
    }
    & & \\ \\

    \theutterance \stepcounter{utterance}  
    & & & \multicolumn{2}{p{0.3\linewidth}}{
        \cellcolor[rgb]{0.9,0.9,0.9}{
            \makecell[{{p{\linewidth}}}]{
                \texttt{\tiny{[GM$|$GM]}}
                \texttt{game\_result = LOSE} \\
            }
        }
    }
    & & \\ \\

\end{supertabular}
}

\end{document}
